% ============================================================
\section{Results}
\label{sec:results}
% ============================================================

\subsection{Predictive Performance}
\label{sec:perf}

Table~\ref{tab:accuracy} reports mean accuracy (up to ten seeds) for each
encoding and the strongest classical baseline per dataset.

\begin{table}[tbp]
\centering
\caption{Mean accuracy (up to ten seeds) for QIE encodings and the best classical
baseline per dataset.  Bold indicates the best QIE result per row.
$\dagger$ denotes all methods near chance (XOR task).}
\label{tab:accuracy}
\small
\begin{tabular}{lrrrr}
\toprule
Dataset & Amplitude & Angle & Basis & Best classical \\
\midrule
Wine                     & 0.628 & \textbf{0.972} & 0.917 & 0.978 (raw linear) \\
Breast Cancer            & 0.779 & \textbf{0.977} & 0.940 & 0.975 (PyTorch MLP) \\
Dry Bean                 & 0.260 & 0.928          & 0.920 & \textbf{0.931} (PyTorch MLP) \\
Credit Fraud             & 0.999 & 0.999          & 0.999 & \textbf{0.999} (MLP) \\
Fashion-MNIST            & 0.821 & 0.823          & 0.819 & \textbf{0.859} (PyTorch MLP) \\
CIFAR-10                 & 0.378 & 0.306          & 0.328 & \textbf{0.457} (RBF-SVM) \\
HIGGS                    & 0.643 & 0.658          & 0.669 & \textbf{0.748} (PyTorch MLP) \\
High-dim parity$\dagger$ & 0.485 & 0.497          & 0.504 & 0.524 (PyTorch MLP) \\
High-rank noise          & \textbf{0.795} & 0.780  & 0.700 & 0.795 (RBF-SVM) \\
Covertype                & 0.587 & 0.725          & 0.724 & \textbf{0.859} (MLP) \\
\bottomrule
\end{tabular}
\end{table}

\subsection{Statistical Analysis}
\label{sec:stats}

Of 30 encoding-dataset comparisons against the best classical baseline, 27 are
statistically significantly worse ($p < 0.05$, paired $t$-test) and 0 are
significantly better.
The three non-significant cases are genuine negligible near-ties: angle on Wine
($d = -0.18$, $p = 0.70$), angle on Breast Cancer ($d = +0.10$, $p = 0.83$),
and amplitude on high-rank noise ($d = -0.03$, $p = 0.96$).
Among the negligible near-ties, the Breast Cancer angle result is directionally
favourable to QIE (mean accuracy 0.977 vs.\ 0.975 for PyTorch MLP), though
the difference is not significant and the effect size is trivially small.
Effect sizes among the significant losses range from small (basis on Breast
Cancer: $d = -1.63$) to catastrophic (amplitude on Dry Bean: $d = -111.9$,
$p < 0.0001$; amplitude on Covertype: $d = -44.0$, $p < 0.0001$).
An important caveat applies to Credit Card Fraud: all three QIE methods are
classified as significantly worse ($p < 0.001$) with large Cohen's $d$
($-4.74$, $-3.97$, $-1.97$), yet the underlying accuracy differences are
$0.0007$, $0.0004$, and $0.0001$ respectively.
These results are statistically significant because near-zero cross-seed
variance inflates the $t$-statistic, not because the performance gap is
practically meaningful.
Figure~\ref{fig:forest} displays the forest plot of all 30 effect sizes.

\begin{figure}[h!]
  \centering
  \includegraphics[width=0.85\textwidth]{figures/forest_plot.png}
  \caption{Forest plot of Cohen's $d$ effect sizes for all 30 encoding-dataset
  comparisons (QIE vs.\ best classical baseline).  Negative values indicate QIE
  is worse.  Error bars are 95\% CIs over five to ten seeds.}
  \label{fig:forest}
\end{figure}
